\section{Síntese dos Resultados}
\label{sec:sintese-resultados}

Os resultados obtidos no benchmark indicam que o Rocket apresentou o melhor desempenho operacional entre os três algoritmos nos dois cenários avaliados. Essa vantagem aparece principalmente na latência média, no percentil 95 da latência e no consumo de memória RAM. O Rocket também apresentou número médio de gerações semelhante ao Morpheus, em razão da ocorrência frequente de parada antecipada.

No cenário i3, o Rocket foi mais rápido tanto na média quanto no p95, além de consumir menos memória que Morpheus e Sunny. O Morpheus foi mais rápido que o Sunny e consumiu muito menos memória, preservando a conclusão da comparação original entre os dois algoritmos. No cenário i5/full, o mesmo comportamento foi mantido, com vantagem operacional do Rocket e vantagem do Morpheus sobre o Sunny na latência média, no p95 e no consumo de memória.

O Sunny apresentou taxa de sucesso total nos dois cenários e teve desempenho pontualmente melhor que o Morpheus na análise do parâmetro taxa de mutação. Também deve ser considerado que sua implementação não se limita à penalização pura, pois utiliza organização interna dos dados, inicialização guiada por restrições e reparo parcial. Ainda assim, no resultado agregado do benchmark, o Sunny consumiu mais memória, executou mais gerações e não obteve redução no erro médio final em relação ao Morpheus e ao Rocket.

A comparação de score favoreceu Morpheus e Rocket, mas com diferença pequena em relação ao Sunny. O erro médio permaneceu igual para os três algoritmos, com valor de 3,00. Assim, não é possível afirmar que Rocket ou Morpheus produziram grades completamente válidas ou qualitativamente superiores em todos os aspectos. O que os dados sustentam com maior clareza é que Rocket e Morpheus produziram resultados equivalentes no erro médio com menor custo operacional que o Sunny, sendo o Rocket o algoritmo com menor custo operacional entre os três.

A inclusão do Rocket revela uma dimensão que a comparação anterior entre Morpheus e Sunny não permitia observar. Como Morpheus e Rocket apresentaram número médio de gerações e frequência de parada antecipada muito semelhantes, mas tempos e consumos de memória bastante diferentes, os resultados sugerem que a eficiência não depende apenas de encerrar o ciclo evolutivo mais cedo. Ela também pode estar associada ao custo por geração, ao modelo de organização dos dados, à avaliação paralela da população, ao reparo local e ao gerenciamento de memória da implementação.

Dessa forma, considerando o \textit{payload}, os parâmetros e o ambiente avaliados, o Rocket mostrou-se mais adequado para o contexto analisado sob a perspectiva operacional. O Morpheus manteve-se como alternativa mais eficiente que o Sunny, enquanto o Sunny demonstrou estabilidade de execução, mas com maior custo de memória e maior número de gerações. Essa conclusão, contudo, está restrita ao protocolo de benchmark executado e não deve ser generalizada como superioridade universal de linguagem, arquitetura ou abordagem evolutiva.
